% Separate post-CP27 candidate. No exact wild-index assertion is inferred.
\subsection{A wild absolute-index cyclic-point restriction}

\begin{lemma}[The number of rational high inverse steps]
\label{lem:B-RD76-point}
\status{Candidate; complete proof supplied, independent audits pending.}
Let \(p\ge5\), and let \(R\) be a complete discrete valuation
ring of mixed characteristic \((0,p)\) with algebraically closed
residue field. Normalize \(v_R(\pi_R)=1\), put
\(a=v_R(p)\) and \(v=v_p(a)\). If a one-dimensional formal
group of special height two has a rational point of exact order
\(p^n\), where \(n\ge2\), then
\[
             a>(p^2-1)p^{2n-3-v}.                 \tag{B-RD76.1}
\]
The right side is interpreted as a positive rational number if
the exponent is negative.
\end{lemma}
\begin{proof}
Use precisely the integral coordinate, coefficient bounds and
Weierstrass polygons in the proof of
Lemma~\ref{lem:B-RD63-point}, from
\cite[Proposition~5.7; (12.3)--(12.5); Proposition~12.10 and proof]
{HopkinsGross1994Bundles}. Those coefficient statements do not
require \(p\nmid a\). Set \(q=p^2\), \(N=q-1\), and
\(b=v_R(u)\), allowing \(b=+\infty\).
If \(b\ge ap/(p+1)\), all nonzero first-level valuations are
\(a/N\), and all subsequent inverse valuations divide by \(q\).
Rationality therefore gives \(a\ge Nq^{n-1}\), which is strictly
stronger than (B-RD76.1).

Otherwise \(b\) is a positive integer and the first-level
valuations are
\[
 \alpha=\frac{a-b}{p-1},\qquad
 \beta=\frac{b}{p(p-1)},\qquad \alpha>\beta>0.
                                                     \tag{B-RD76.2}
\]
For an input of valuation \(0<Z<a\), the exact inverse polygon
gives
\[
 \begin{cases}
 Z/q,& b\ge Z(p-1)/p,\\
 (Z-b)/p\text{ with multiplicity }p,
 \quad\beta\text{ with multiplicity }q-p,
      &b<Z(p-1)/p.
 \end{cases}                                         \tag{B-RD76.3}
\]
Equality belongs to the uniform first case. Each new valuation
is smaller than its input; all input valuations are below \(a\).
After either a uniform step or a low step to \(\beta\), every
remaining step is uniform, as shown in that same proof.

Let \(P\) be the asserted point. Every point in its successive
multiplication chain is rational, so every valuation is a
positive integer. If the first-level point has valuation
\(\beta\), then
\[
 b\ge p(p-1)q^{n-1},\qquad
 a>\frac{p+1}{p}b\ge Nq^{n-1},
\]
again stronger than required. It remains that the chain starts
at \(Z_0=\alpha\). After \(j\) consecutive high inverse steps,
the exact recurrence in (B-RD76.3) gives
\[
 Z_j=\frac{a/p^j-b}{p-1},\qquad
 a=p^j\bigl(b+(p-1)Z_j\bigr).
                                                     \tag{B-RD76.4}
\]
Integrality of \(Z_j\) and \(b\) forces \(p^j\mid a\), hence
\(j\le v\). This is the only new constraint replacing the
prime-to-\(p\) hypothesis in B63.

Suppose first that the first non-high step is uniform, after
\(j\le\min(v,n-2)\) high steps. There are
\(r=n-1-j\ge1\) uniform inverse steps left. Rationality of
the terminal valuation implies \(Z_j\ge q^r\), and the
uniform inequality implies \(b\ge Z_j(p-1)/p\). Therefore
\[
 a=p^j\bigl(b+(p-1)Z_j\bigr)
    \ge Np^{j-1}Z_j
    \ge Np^{2n-3-j}
    \ge Np^{2n-3-v}.                              \tag{B-RD76.5}
\]
If \(j<v\), the last inequality is strict. If \(j=v\),
equality throughout would give
\(v_p(a)=2n-3-v>v\), because \(v=j\le n-2\).
This contradicts the definition of \(v\). The bound is thus
strict in this case as well.

If the first non-high step is low, after \(j\) high steps,
then the remaining \(n-2-j\) inverse steps divide \(\beta\)
by \(q\). Integrality and the strict high-case inequality give
\[
 \begin{gathered}
 b\ge p(p-1)q^{n-2-j},\qquad
 a/p^j>(p+1)b,\\
 a>Np^{2n-3-j}\ge Np^{2n-3-v}.
 \end{gathered}                                    \tag{B-RD76.6}
\]
Finally, if all \(n-1\) inverse steps are high, then
\(v\ge n-1\) and \(Z_{n-1}\ge1\). Since \(b\ge1\),
\[
 a=p^{n-1}\bigl(b+(p-1)Z_{n-1}\bigr)
       \ge p^n>Np^{n-2}\ge Np^{2n-3-v}.
                                                     \tag{B-RD76.7}
\]
This exhausts every possible chain, including each equality
boundary of the Newton polygon.
\end{proof}

\begin{proposition}[A mixed comparison at wild absolute index]
\label{prop:B-RD76-mixed}
\status{Candidate; complete proof supplied, independent audits pending.}
Let \(K\) be complete discretely valued of mixed characteristic
\((0,p)\), with perfect residue field and \(p\ge5\). Put
\(e=v_K(p)\) and \(v=v_p(e)\). Let \(M\) be a height-two
BT group which acquires one-dimensional supersingular good
reduction, and let \(H\) have a two-step potentially good
height-one filtration of total dimension one. For \(n\ge2\),
\[
 M[p^n]\simeq H[p^n]\quad\Longrightarrow\quad
                     e>p^{2n-3-v}.                 \tag{B-RD76.8}
\]
\end{proposition}
\begin{proof}
Pass to the completed maximal unramified extension, preserving
\(e\), \(v\), and the comparison. The quotient-and-finite-twist
argument of Proposition~\ref{prop:B-RD63-mixed}, whose complete
construction is in Proposition~\ref{prop:B-RD49-mixed}, produces
a potentially supersingular height-two group with a rational
point of exact order \(p^n\). None of those finite-flat closure,
actual quotient, height-one character, or finite twist steps
assumes \(p\nmid e\).

Lemma~\ref{lem:B-RD49-inertia} supplies a good formal model
over a tame extension of degree \(m\mid N=p^2-1\), for every
absolute index \(e\). The rational point extends to its finite
flat torsion model and lies in the formal group, because the
special fibre is connected. Here \(a=me\) and \(v_p(a)=v\).
The preceding lemma gives \(me>Np^{2n-3-v}\). Since
\(m\le N\), (B-RD76.8) follows.
\end{proof}

\begin{lemma}[Dimension detection with a wild absolute index]
\label{lem:B-RD76-determinant}
\status{Candidate; complete proof supplied, independent audits pending.}
In the height-two monodromy setting of B.4, assume \(p\ge5\),
\(\mu_p\subset K\), perfect residue field, and write
\(v=v_p(e)\). If \(n\ge v+2\), an isomorphism of actual
\(p^n\)-torsion determines the total dimension of the potentially
good graded pieces and their finite determinant character.
\end{lemma}
\begin{proof}
Pass again to the completed maximal unramified base. The generic
determinant on inertia is \(\chi_p^d\epsilon\), where
\(d\in\{0,1,2\}\) and \(\epsilon\) is
\(\mu_{p-1}\)-valued, by the determinant proof in B49.
Since \(\mu_p\subset K\), the extension \(K(\mu_{p^n})/K\)
is cyclic of degree \(p^r\). Its residue field is unchanged.
In its normalized valuation,
\[
 v_{K(\mu_{p^n})}(\zeta_{p^n}-1)
     =\frac{ep^r}{p^{n-1}(p-1)}\in\mathbf Z_{>0},
 \qquad r\ge n-1-v.                                \tag{B-RD76.9}
\]
Thus the cyclotomic character modulo \(p^n\) has a nontrivial
\(p\)-part when \(n\ge v+2\). Equality of determinants forces
\(p\mid d_1-d_2\), whereas \(|d_1-d_2|\le2<p\), so
\(d_1=d_2\). Their remaining finite characters agree, since
reduction modulo \(p\) is injective on \(\mu_{p-1}\).
\end{proof}

\begin{corollary}[The remaining branch at \(p=5,e=20\)]
\label{cor:B-RD76}
\status{Candidate; complete proof supplied, independent audits pending.}
Under the full B.4 height-two hypotheses with \(p\ge5\),
\(\mu_p\subset K\), and perfect residue field, suppose
\[
 n\ge\max\{2,v_p(e)+2\},\qquad
 e\le p^{2n-3-v_p(e)}.                              \tag{B-RD76.10}
\]
Any failure of special-character determination by actual
\(p^n\)-torsion must compare two whole potentially
supersingular height-two groups. In particular, at \(p=5,e=20\),
any failure at actual \(125\)-torsion is confined to that branch;
the mixed whole/filtered branch is impossible.
\end{corollary}
\begin{proof}
The determinant lemma fixes the total dimension. Dimension zero
is determined by the Brauer lift of semisimplified residual
inertia; dimension two reduces to it by Cartier duality, exactly
as in B49. In dimension one, each allowed group is either whole
potentially supersingular or filtered by two potentially good
height-one pieces. The characters of two filtered groups are
determined residually because \(\chi_p\equiv1\pmod p\).
The mixed possibility is excluded by (B-RD76.8) and the stated
inequality. At \(p=5,e=20\) we have \(v_p(e)=1\), and
\(n=3\) satisfies \(20\le5^2\) and \(n=v_p(e)+2\).
\end{proof}

\begin{remark}[No wild least-index assertion]
\label{rem:B-RD76-scope}
\status{Candidate scope; no whole-formal recovery or sharpness claimed.}
For \(v_p(e)=0\), the mixed necessary inequality is exactly B63.
For positive \(v_p(e)\), it is only a necessary inequality; no
matching existence construction at its new boundary is asserted.
The new \(p=5,e=20\) conclusion eliminates a comparison class
and does not prove that \(125\)-torsion determines every allowed
character. B63 already supplies failure at actual \(25\)-torsion
for this index, so its least manuscript index is at least two.
The whole-potentially-supersingular \(125\)-comparison remains
open. All claims here use the actual torsion level, not a rational
isogeny class or a count of numerical Newton possibilities.
\end{remark}
